\documentclass{article}

\usepackage{amsmath,amssymb}
\usepackage{xurl}
\usepackage[hidelinks]{hyperref}
\setlength{\emergencystretch}{2em}

% Shared commands for notes in docs/paper-gaps/.

\DeclareUnicodeCharacter{2080}{\ifmmode{}_{0}\else\textsubscript{0}\fi}
\DeclareUnicodeCharacter{2081}{\ifmmode{}_{1}\else\textsubscript{1}\fi}
\DeclareUnicodeCharacter{2082}{\ifmmode{}_{2}\else\textsubscript{2}\fi}
\DeclareUnicodeCharacter{2099}{\ifmmode{}_{n}\else\textsubscript{n}\fi}

\newcommand{\Lean}{\textsc{Lean}}
\ExplSyntaxOn
\NewDocumentCommand{\leanid}{m}{
  \ifmmode
    \textsf{\detokenize{#1}}
  \else
    \group_begin:
    \urlstyle{sf}
    \tl_set:Nn \l_tmpa_tl {#1}
    \tl_replace_all:Nnn \l_tmpa_tl {\_} {_}
    \exp_args:NV \path \l_tmpa_tl
    \group_end:
  \fi
}
\ExplSyntaxOff
\providecommand{\tr}{\operatorname{tr}}
\newcommand{\tnleanIssueBase}{https://github.com/LionSR/TNLean/issues/}
\newcommand{\tnleanPrBase}{https://github.com/LionSR/TNLean/pull/}
\newcommand{\tnleanDocBase}{https://sirui-lu.com/TNLean/docs/}
\newcommand{\ghissue}[1]{\href{\tnleanIssueBase#1}{GitHub issue \##1}}
\newcommand{\ghpr}[1]{\href{\tnleanPrBase#1}{PR \##1}}
\newcommand{\doclink}[1]{\href{\tnleanDocBase#1.html}{\path{#1}}}

% Dirac notation and the identity operator, as in blueprint/src/macros/common.tex.
% \providecommand keeps note-local definitions authoritative.
\usepackage{dsfont}
\providecommand{\ket}[1]{|#1\rangle}
\providecommand{\bra}[1]{\langle #1|}
\providecommand{\braket}[2]{\langle #1 | #2 \rangle}
\providecommand{\ketbra}[2]{|#1\rangle\!\langle #2|}
\providecommand{\Id}{\mathds{1}}
\providecommand{\one}{\mathds{1}}
\providecommand{\C}{\mathbb{C}}
\providecommand{\MN}[1]{M_{#1}(\C)}
\providecommand{\MD}{M_D(\C)}

% External referencing for paper sources.  The build helper writes sanitized
% auxiliary files under build/paper-aux/ and excludes duplicated source labels.
\usepackage{xr}

% Import a generated paper auxiliary file with a caller-supplied label prefix.
% The candidate paths support builds from docs/paper-gaps/, the repository root,
% and build/paper-aux/ itself.
\newcommand{\importpaperaux}[2]{%
  \IfFileExists{../../build/paper-aux/#2.aux}{%
    \externaldocument[#1]{../../build/paper-aux/#2}%
  }{%
    \IfFileExists{build/paper-aux/#2.aux}{%
      \externaldocument[#1]{build/paper-aux/#2}%
    }{%
      \IfFileExists{#2.aux}{%
        \externaldocument[#1]{#2}%
      }{}%
    }%
  }%
}

% General external reference macros
\newcommand{\paperref}[2]{\ref{#1-#2}}
\newcommand{\papereqref}[2]{(\ref{#1-#2})}

% CPSV16 source (1606.00608 / MPDO-22-12-17-2.tex) external document and shortcuts
\importpaperaux{cpsv16-}{1606.00608/MPDO-22-12-17-2-tnlean}
\newcommand{\cpsvref}[1]{\paperref{cpsv16}{#1}}
\newcommand{\cpsveqref}[1]{\papereqref{cpsv16}{#1}}

% Verdict marker: \gapnote{<kind>}{<status>}. Kind is the policy
% classification (clarification, local-correction, scope-restriction,
% unfaithful, false-source, open-gap); status is open, wip (elimination
% underway), resolved, or historical. Machine-read by the site index;
% prints a verdict line.
\newcommand{\gapnote}[2]{\par\noindent\textbf{Verdict:} #1 (#2).\par}


\title{The Phase Index in the Weighted Ces\`aro Formula, Wolf Equation~(6.15)}
\author{}
\date{2026-08-16}

\begin{document}
\maketitle

\section{The source sentence}

After expressing $T_\infty$ as a Ces\`aro mean in his Equation~(6.14), Wolf
writes, for a trace-preserving positive linear map $T:\MN{D}\to\MN{D}$,
\begin{align}
  T_\phi
  & = \lim_{N\to\infty}\frac{1}{N}\sum_{n=1}^N
  \sum_{k:|\lambda_k|=1} (\overline{\lambda_k}T)^n.
  \tag{Wolf eq.~6.15}
\end{align}
The local source is
\path{Notes/WolfNoteTexSource/ch06_spectral_properties.tex},
lines~258--264 \cite{Wolf2012Quantum}. Throughout the chapter the index $k$
runs over the Jordan blocks of the superoperator $\widehat T$: by
Equations~(6.4)--(6.5) of the same source (lines~111--139),
$\widehat T=X(\bigoplus_{k=1}^K J_k(\lambda_k))X^{-1}$, and ``the number of
Jordan blocks with eigenvalue $\lambda$ is the geometric multiplicity of
$\lambda$'' (lines~140--142); several blocks may carry the same eigenvalue.

\section{Correction forced by the multiplicities}

The summand $(\overline{\lambda_k}T)^n$ depends on the eigenvalue $\lambda_k$
only, not on the block carrying it, so the literal block-indexed inner sum
counts each peripheral eigenvalue with its geometric multiplicity. Take the
identity channel $T=\one$ on $\MN{D}$ with $D>1$: every matrix is an eigenvector
for the eigenvalue $1$, the Jordan decomposition of $\widehat T$ consists of
$D^2$ one-dimensional blocks all carrying $\lambda_k=1$, and the literal
right-hand side of Equation~(6.15) is
\[
  \lim_{N\to\infty}\frac{1}{N}\sum_{n=1}^N D^2\,\one
  = D^2\,\one
  \neq T_\phi=\one.
\]
So read literally, with $k$ a block index, the equation is false. The reading
under which it is true --- plainly the intended one, since the same chapter
characterizes $T_\phi$ as the projection onto the span of the peripheral
eigenvectors - is that $k$ runs over the \emph{distinct} peripheral
eigenvalues,
\[
  T_\phi
  = \lim_{N\to\infty}\frac{1}{N}\sum_{n=1}^N
  \sum_{\substack{\lambda\ \text{eigenvalue of }T\\ |\lambda|=1}}
  (\bar\lambda T)^n ,
\]
each phase counted once. This is consistent with the rest of the chapter:
peripheral Jordan blocks of a trace-preserving positive map are
one-dimensional (source, proposition on lines~181--186), so the peripheral
subspace is the direct sum of the eigenspaces of the distinct peripheral
eigenvalues, and on the eigenspace of $\lambda$ the summand indexed by $\mu$
acts as the scalar $(\bar\mu\lambda)^n$, whose Ces\`aro mean is
$\delta_{\mu\lambda}$; a single summand per distinct eigenvalue therefore
suffices to recover the projection. The neighbouring sums
$\widehat T_\infty$, $\widehat T_\phi$, $\widehat T_\varphi$ of
Equations~(6.11)--(6.13) are indexed by blocks as well, but they involve the
block projections $P_k$ and are correct as printed; only Equation~(6.15),
whose summand lost the block dependence, requires the deduplicated index.

The formal development averages over the set of eigenvalues of modulus one, so
each distinct phase appears once.\footnote{Formal declarations:
  \leanid{Module.End.weightedCesaroMean} for the average, and the convergence
  statements
  \leanid{IsPositiveMap.tendsto_weightedCesaroMean_peripheralProjection},
  \leanid{IsPositiveMap.tendsto_endEquiv_weightedCesaroMean_peripheralProjection},
  \leanid{IsPositiveMap.tendsto_weightedCesaroMean_toFinset_peripheralProjection},
  and
  \leanid{IsPositiveMap.tendsto_endEquiv_weightedCesaroMean_toFinset_peripheralProjection}
  in \doclink{TNLean/Channel/Peripheral/WeightedCesaro}.}
This correction changes no hypothesis of Equation~(6.15); it adjusts the
multiplicity convention of the inner summation index.

\section{Verdict}

This is a local correction of the summation index in the source equation: $k$
is reinterpreted from a Jordan-block index to an index of the distinct
peripheral eigenvalues. It does not alter the hypotheses or the conclusion of
Equation~(6.15), which is false under the literal block-indexed reading, as
the identity channel shows.

\bibliographystyle{alpha}
\bibliography{references}

\end{document}